\documentclass[11pt,a4paper]{article}
\usepackage[utf8]{inputenc}
\usepackage[T1]{fontenc}
\usepackage[spanish]{babel}
\usepackage{hyperref}
\usepackage{geometry}
\geometry{margin=2.5cm}
\usepackage{parskip}

\title{113 la tesis no se sostiene en un centro de datos}
\author{Alejandro Valenzuela Rivera\\\small Universidad Autónoma de Nuevo León}
\date{\today}
\begin{document}
\maketitle

\begin{quote}
\textbf{Sesgo declarado.} Este trabajo se sostiene desde el sesgo declarado de su autor, y no finge neutralidad. Lo que afirma se puede tumbar, y ese es el punto: el falsificador está escrito, no escondido.
\end{quote}

\section*{Resumen}
Se sostiene en una consola de mano de 16 GB que llama a una API. Ese es el punto: si corre aqui, corre en casi cualquier cosa. Y si el metodo no depende de la maquina, el que venga puede correrlo con un modelo local modesto y mas paciencia.

\section{Por qué se sostiene}
Y lo que cambia entre una API y un modelo local no es el metodo ni la calidad del piso: es el tiempo y el costo por resultado, y los dos se miden con STATS. Lo que no cambia es que el que lo corre pueda contestar.

\section{De dónde viene}
La autoridad, 2026-10-26.

\section{El caso}
Ninguna: es la tesis aplicada a su autor.

\section{Método}
Una tesis que se sostiene en lo que una persona puede correr en su maquina.

\section{Límites, y qué lo tumba}
Una tesis que solo se demuestra con recursos que nadie mas tiene.

\section{Lo que queda abierto}
\textquestiondown{}Que se demuestra con una maquina modesta?

\section*{Cómo se leyó esta tesis}
Esto se escribió en el repositorio man-made-empathy, del que este candado es una pieza: candados/06 del caso/113 la tesis no se sostiene en un centro de datos.md. La máquina que lo compuso lo hizo desde este mismo archivo, y no agregó nada que no estuviera escrito. Si algo de esto es falso, se dice en el mismo sitio, y se corrige con fecha.

\end{document}
